\section{Finite observations and amplification}\label{sec:observations}

We first justify the degree and conditioning identities used by every
construction. Throughout, \(F\) is a finite-valued observation,
\(T_F=\E[H_N\mid F]\) is its centered score, and \(D>0\) is a cell
degree bound. The bound \(D\) may be real; it need not equal the actual
integer degree of any indicator.

Every real function of \(F\) has degree at most \(D\), by expansion in
its finitely many value indicators. Unused labels have zero indicators.
If distinct observation labels produce the same numerical score, their
indicators add, so score collisions preserve the same bound. A function of \(r\) observations on
disjoint blocks, with bounds \(D_1,\ldots,D_r\), has degree at most
\(\sum_iD_i\). Indeed, it expands in products of their indicators, and
the degree of a product is at most the sum of the degrees: multiplication
of Fourier characters takes symmetric differences of their index sets.

Let \(F_1,\ldots,F_r\) be independent observations on disjoint bit
blocks, with scores \(T_1,\ldots,T_r\), and let \(H_{\mathrm{tot}}\)
be the total bit sum. Conditioning on their tuple gives
\(\E[H_{\mathrm{tot}}\mid F_1,\ldots,F_r]=\sum_iT_i\).
For any observation \(Q\) that is a function of this tuple, the tower
property therefore gives
\[
 \E[H_{\mathrm{tot}}\mid Q]=\E\!\left[\sum_iT_i\,\middle|\,Q\right].
\]
These are instances of the conditional-expectation projection identities;
see~\cite[Theorems~4.1.13 and 4.1.15]{Durrett}. In particular, for
\(v=v(F)>0\), the numerical observation
\begin{equation}\label{eq:averaging}
 Z_M=\frac{T_1+\cdots+T_M}{\sqrt{Mv}}
\end{equation}
formed from \(M\) independent copies has cell degree bound \(MD\), and
its score is exactly \(\sqrt{Mv}\,Z_M\).
The classical i.i.d.\ central limit theorem~\cite[Theorem~3.4.1]{Durrett} gives
\(Z_M\Rightarrow G\), where \(G\) is standard normal. For each fixed
initial \(F\), all moments are finite, and
\begin{equation}\label{eq:fourth}
 \E Z_M^4=\frac{\E T_F^4}{Mv^2}+3\frac{M-1}{M}.
\end{equation}
Thus the fourth moments are bounded uniformly in \(M\) for this fixed
observation. The bound may change when a later amplification stage
uses a different observation.
All citations to probability results refer to the January 11, 2019 author
version of Durrett~\cite{Durrett}.

\begin{lemma}[Transfer of predictor errors]\label{lem:transfer}
Let \(Z_j\) be centered, variance-one random variables converging in
distribution to a standard normal, with \(\sup_j\E Z_j^4<\infty\).
For fixed \(r\), let \(\mathbf Z_j\) consist of \(r\) independent copies
of \(Z_j\), and let \(\mathbf G\) consist of \(r\) independent standard
normals. If \(g:\R^r\to\R\) is measurable, its discontinuities have
Gaussian measure zero, and
\[
 |g(z)|\le K\left(1+\sum_{i=1}^r z_i^2\right)
\]
for some fixed \(K\), then \(\E g(\mathbf Z_j)\to\E g(\mathbf G)\).
\end{lemma}

\begin{proof}
Independence gives joint weak convergence. Clipping \(g\) to a bounded
range gives convergence of expectations by the null-discontinuity
criterion for weak convergence~\cite[Theorem~3.10.1(vi)]{Durrett}.
The fourth-moment bound and fixed \(r\)
give \(\sup_j\E[g(\mathbf Z_j)^2]<\infty\), so the clipping errors in
first moment tend to zero uniformly; this is the higher-moment
criterion for uniform integrability~\cite[Theorem~4.6.2]{Durrett}.
The Gaussian error does also.
\end{proof}

\begin{lemma}[Boolean readout]\label{lem:readout}
Suppose finite observations have positive cell degree bounds \(D\) and
arbitrarily large ratios \(v(F)/D\). Then Theorem~\ref{thm:main} holds.
\end{lemma}

\begin{proof}
Fix \(C>0\) and choose \(F\) with \(v(F)>4C^2D\).
For \(M\) independent copies, define
\(f=\operatorname{sign}(T_1+\cdots+T_M)\), with sign \(1\) at zero.
Its degree is at most \(MD\). Conditioning on the tuple of observations,
\[
 \sum_i\widehat f(\{i\})
 =\E[fH_{MN}]
 =\E|T_1+\cdots+T_M|
 =\sqrt{Mv(F)}\,\E|Z_M|.
\]
The central limit theorem and the uniform second-moment bound imply
\(\E|Z_M|\to\E|G|=\sqrt{2/\pi}>1/2\).
For a sufficiently large finite \(M\), the displayed sum exceeds
\(C\sqrt{MD}\ge C\sqrt{\deg(f)}\). It is positive, so \(f\) is
nonconstant.
\end{proof}

\begin{proposition}[Amplification criterion]\label{prop:amplification}
Let \(1\le d<r\) be integers, and let \(\mathcal W\) be a measurable
observation rule on \(r\) real inputs. Suppose:
\begin{enumerate}
\item On any product of finite input alphabets, the indicator of each
output value is a finite linear combination of functions each depending
on at most \(d\) input coordinates.
\item There is an \(\eta>0\) such that, whenever centered variance-one
laws \(Z_j\) converge weakly to a standard normal with uniformly bounded
fourth moments, independent copies satisfy, for all sufficiently large \(j\),
\[
 \Var\!\left(\E\!\left[\sum_{i=1}^r Z_{j,i}\,\middle|\,
                \mathcal W(\mathbf Z_j)\right]\right)\ge d+\eta.
\]
\end{enumerate}
Then finite observations have arbitrarily large retained
variance-to-cell-degree ratios, and Theorem~\ref{thm:main} follows.
\end{proposition}

\begin{proof}
Fix a finite observation \(F\) with \(v=v(F)>0\) and degree bound \(D>0\).
Form \(r\) independent blocks of \(M\) copies of \(F\), and apply
\(\mathcal W\) to their normalized scores from \eqref{eq:averaging}.
The new observation \(F'\) has finite range. Its cell degree is at most
\(dMD\) by the first assumption and expansion in value indicators.
The second assumption applies by \eqref{eq:fourth}. For a sufficiently
large finite \(M\), the tower property therefore gives
\[
 v(F')=Mv\,\Var\!\left(\E\!\left[\sum_{i=1}^rZ_{M,i}\mid F'\right]\right)
 \ge Mv(d+\eta).
\]
The ratio increases by at least the fixed factor \(1+\eta/d>1\).
Start with the observation of one sign, for which \(v=D=1\), and
iterate. Each stage is finite and has its own sufficiently large finite
copy count. A finite number of stages exceeds any prescribed ratio.
Apply Lemma~\ref{lem:readout}.
\end{proof}
